% Completed ordinary calculation; independently audited by root and fhi.
\subsection*{D4o: the complete ordinary Igusa deformation calculation}
\textbf{Status: Proved.} The calculation was independently audited by root and fhi, including the opened Katz and Mumford--Fogarty sources. It covers ordinary first factors; the supersingular case is completed separately in D4s.

Let $k$ be algebraically closed of characteristic two, let $E_0,F_0$ be ordinary
elliptic curves, and let $P\in F_0[2](k)$ be nonzero. Put
\[
 A_0=E_0\times F_0,\quad \sigma_0=(-1,1),\quad
 \tau_0=t_{(0,P)}\sigma_0,\quad X_0=A_0/\langle\tau_0\rangle.
\]
The base category consists of local Artinian $W(k)$-algebras with residue
field $k$, and all deformations are marked on the closed fibre.

\paragraph{Step 1: lifting the cover and choosing its origin.}
The finite etale cover $A_0\to X_0$, together with its deck involution,
lifts uniquely to every deformation $X/R$, by invariance of the etale site
under a nilpotent immersion. This equivalence also applies to marked
isomorphisms: specify the identity lift on $A_0$ to remove the alternative
deck transformation. The lifted cover $A/R$ is smooth proper. A section
lifting $0\in A_0(k)$ exists by smoothness. It gives a unique abelian-scheme
structure, by Mumford--Fogarty, Proposition~6.15, pp.~124--126
\cite{DGJ_MumfordFogarty1982}, applied successively to small extensions.
This proposition assumes properness, not projectivity.

Write the deck map uniquely as $\tau=t_c\sigma$, where
$c=\tau(0)$ and $\sigma$ fixes the origin. Rigidity makes $\sigma$ a group
automorphism lifting $\sigma_0$, and $\sigma^2=1$. Thus the only additional
equation for the affine map is $(1+\sigma)c=0$.

Conversely every such triple $(A,\sigma,c)$ gives a free involution: its
fixed locus is closed in the proper scheme $A$ and has empty special
fibre, hence is empty. Its quotient is a smooth proper deformation of
$X_0$. At each Artinian stage $A$ is projective: a sufficiently high
power of an ample bundle on $A_0$ lifts through successive small
extensions, because their obstruction groups are killed by two.
An invariant ample bundle is obtained by tensoring a bundle with its
translate, so the quotient is projective too.

\paragraph{Step 2: Serre--Tate parameters and exact point coordinates.}
Choose bases $e_1,e_2$ of $T_2A_0(k)$ and $e_1^\vee,e_2^\vee$ of
$T_2A_0^\vee(k)$, with $e_2/2=P$. Use Katz's convention
\[
 q_{ij}=q(A/R;e_i,e_j^\vee)\in1+\mathfrak m_R.
\]
Katz, Theorem~2.1, pp.~148--149, identifies deformations with this matrix,
identifies duality with transposition, and gives the homomorphism-lifting
criterion. Therefore $\sigma_0$ lifts exactly when
\[
 q_{12}^2=q_{21}^2=1.                                      \tag{I1}
\]

Here is an $R$-point interpretation that avoids any unstated surjectivity
of a torsion subgroup. Choose a compatible system
$P_n\in F_0[2^{n+1}](k)$ with $2^nP_n=P$. For $n$ sufficiently large,
$[2^n]$ kills the formal group $\widehat A(\mathfrak m_R)$.
Choose an arbitrary lift $\widetilde P_n\in A(R)$ and put
\[
 c_0=[2^n]\widetilde P_n.
\]
It is independent of the lift and of sufficiently large $n$.
Indeed changes of lift lie in the killed formal group; replacing
$n$ by $n+1$ gives the same result because
$2\widetilde P_{n+1}-\widetilde P_n$ lies in that formal group.
It is functorial in $R$, reduces to $(0,P)$, and is fixed by $\sigma$,
because $\sigma_0P_n=P_n$. By Katz's explicit recovery of the
pushout homomorphism, pp.~150--152,
\[
 [2]c_0=(q_{21},q_{22})\in
 \widehat A(\mathfrak m_R)\simeq(1+\mathfrak m_R)^2.        \tag{I2}
\]
The isomorphism on the right is the one furnished by the chosen dual
Tate basis. It also identifies $\sigma$ on the formal group with
$(z_1,z_2)\mapsto(z_1^{-1},z_2)$.

Every lift $c$ of $(0,P)$ is uniquely $c_0+z$, with
$z=(z_1,z_2)\in(1+\mathfrak m_R)^2$. Consequently
\[
 (1+\sigma)c=(q_{21},q_{22}z_2^2).
\]
The affine involution condition is therefore exactly
\[
 q_{21}=1,\qquad q_{22}z_2^2=1.                           \tag{I3}
\]
Together with (I1), the functor of triples is represented by
\[
 R_{\rm tr}=W(k)[[t,s,u,v]]/(s^2+2s),                     \tag{I4}
\]
where
\[
 q_{11}=1+t,\quad q_{12}=1+s,\quad q_{21}=1,\quad
 q_{22}=(1+v)^{-2},\qquad z_1=1+u,\ z_2=1+v.
\]
This inverse in $q_{22}$ fixes Katz's convention; replacing $q$ by
$q^{-1}$ yields the earlier equivalent coordinate draft.

\paragraph{Step 3: the ring is a miniversal hull for the unpointed surface.}
Forgetting the chosen origin gives a formally smooth transformation
from the triple functor to $\operatorname{Def}_{X_0}$: over a small
extension, lift the cover uniquely and lift its prescribed origin by
smoothness. At first order this transformation is injective as well
as surjective. In fact any marked isomorphism of the unpointed covers
is a translation followed by a marked abelian isomorphism; rigidity
forces the latter to preserve the same $q$ parameters. Translation
by a formal point $d$ changes $c$ by $(1-\sigma)d$. In formal torus
coordinates this changes $z_1$ by a square and leaves $z_2$ fixed.
On $k[\epsilon]/(\epsilon^2)$ every principal unit has square one.
Thus all four tangent parameters in (I4) remain distinct. This proves
formal smoothness and tangent bijectivity, the defining hull property.
It also explains why the ring represents triples but need not represent
the surface deformation functor: higher-order changes of origin can
identify different triples.

The compatible formal family can be algebraized over
$\operatorname{Spec}R_{\rm tr}$. Let $\lambda_0$ be the product principal
polarization on $A_0$. The homomorphism $2\lambda_0$ lifts through all
stages, since its Serre--Tate condition is $q_{ij}^2=q_{ji}^2$, which
holds by (I1) and (I3). The canonical diagonal pullback of the
Poincare bundle attached to this symmetric homomorphism gives a
compatible relatively ample line bundle whose special-fibre class is
$4\lambda_0$. Formal existence algebraizes the abelian scheme and
the involution, its translation point, and the quotient. The hull is
unpolarized; the auxiliary polarization is used only for algebraization.

\paragraph{Step 4: the fixed dual group is flat.}
Let $\mathcal E^\vee,\mathcal F^\vee$ denote the elliptic deformations
with Serre--Tate coordinates $q_{11},q_{22}$. The condition $q_{21}=1$
and the transposition rule lift both the original inclusion and projection
to an exact sequence of abelian schemes
\[
 0\longrightarrow\mathcal E^\vee\longrightarrow A^\vee
 \stackrel{\rho}{\longrightarrow}\mathcal F^\vee
 \longrightarrow0.                                      \tag{I5}
\]
For precision, the inclusion condition is the $(1,2)$ entry of the
dual matrix, namely $q_{21}=1$; the projection has the same condition.
The compositions are zero by rigidity. The lifted inclusion is a closed
immersion and the projection is smooth with connected elliptic kernel,
as can be checked on the unique geometric fibre and by the smoothness
criterion; their kernel is therefore the stated elliptic scheme.
Rigidity shows that $\sigma^\vee$ acts by $-1$ on the kernel and by
$+1$ on the quotient.

Put $G=\ker(1-\sigma^\vee:A^\vee\to A^\vee)$.
The restricted map $G\to\mathcal F^\vee$ is an epimorphism of fppf
sheaves. To lift $b\in\mathcal F^\vee(T)$, first lift it fppf locally
to $a\in A^\vee$. The element $(1-\sigma^\vee)a$ lies in
$\mathcal E^\vee$. Divide it by two fppf locally in that elliptic
scheme and subtract the chosen half from $a$. The result is fixed.
Its kernel is exactly $\mathcal E^\vee[2]$. Hence
\[
 0\longrightarrow\mathcal E^\vee[2]\longrightarrow G
 \longrightarrow\mathcal F^\vee\longrightarrow0          \tag{I6}
\]
is fppf exact, and $G\to\mathcal F^\vee$ is a finite locally free
torsor of degree four. In particular $G/R$ is flat.

\paragraph{Step 5: the relative Picard scheme is flat.}
For the finite etale double cover $\pi:A\to X$, translation acts
trivially on $\operatorname{Pic}^0_{A/R}=A^\vee$. Pullback gives an
exact sequence of fppf sheaves
\[
 0\longrightarrow\mu_2\longrightarrow
 \operatorname{Pic}^{\tau}_{X/R}\stackrel{\pi^*}{\longrightarrow}
 G\longrightarrow0.                                    \tag{I7}
\]
Here are the descent details. An invariant relative class is represented
fppf locally by a line bundle $L$. Invariance of its relative class gives an isomorphism up to a base line bundle; trivialize that line bundle fppf locally to obtain
$\alpha:\tau^*L\simeq L$. The square of this isomorphism is a base
unit, since $\Gamma(A_T,\mathcal O)=\Gamma(T,\mathcal O)$ universally.
Adjoin a square root of that unit fppf locally and rescale $\alpha$;
it then obeys the cocycle condition, so $L$ descends to $X$.
The descended class $D$ satisfies
$D^{\otimes2}=\operatorname{Nm}_{A/X}(L)$ in the relative Picard
functor. Fibrewise the norm of a class in $A^\vee$ belongs to the
neutral Picard component, because the norm maps the connected
abelian variety $A^\vee$ to that component. Thus $D$ belongs to
$\operatorname{Pic}^{\tau}$.
Conversely a numerically trivial class pulls back into $A^\vee$, as
the Neron--Severi group of an abelian variety is torsion free.
The choices of linearization on the trivial bundle form precisely
$\operatorname{Hom}(\mathbb Z/2,\mathbb G_m)=\mu_2$, proving the
kernel assertion. These arguments work over arbitrary base changes
$T\to\operatorname{Spec}R$, not only at geometric points.

Thus $\operatorname{Pic}^{\tau}_{X/R}\to G$ is a finite locally free
$\mu_2$-torsor of degree two. Composed with (I6), it is finite locally
free of degree eight over $\mathcal F^\vee$, and is flat over $R$.
In particular the universal first-order ordinary Igusa deformation has
flat $\operatorname{Pic}^{\tau}$.
This refutes the literal conclusion of Chai's Theorem~2.4(iii),
while leaving its line-bundle obstruction and differing tangent-space
computations entirely compatible with flatness.

\paragraph{Opened source locators and independent review.}
Katz \cite{DGJ_Katz1981}, pp.~148--152, contains the coordinate,
duality, lift criterion, and pushout recovery used above. The finite
etale lifting statement is Stacks, Theorem~59.45.2, tag04DY,
\url{https://stacks.math.columbia.edu/tag/04DY}; quotient existence is
Lemma~66.14.2, tag07S7,
\url{https://stacks.math.columbia.edu/tag/07S7}. Both were opened.
Mumford--Fogarty, Proposition~6.15, pp.~124--126, was opened through
the primary-book OCR mirror listed in the bibliography.
The full seven-page Chai v2 was opened in web retrieval, and its final
flatness inference has no additional qualification beyond the tangent
comparison already tested by the explicit torsor example.

No classification beyond ordinary $E_0$ is claimed in this proof. Root and fhi independently checked (I2)--(I7), the origin quotient, and algebraization; fhi's audit is recorded separately. The flatness assertion for every Artinian ordinary deformation uses Steps1,2,4,5 alone and does not depend on miniversality.

\paragraph{Two concrete first-order checks.}
Over $k[\epsilon]/(\epsilon^2)$ the parameters $q_{12}=1+\epsilon$, all other $q_{ij}=1$, and $z_1=z_2=1$ give an actual nonproduct direction of the surface deformation, with flat torsion Picard scheme by (I6)--(I7). There is also a completely explicit product direction: take $E_\epsilon:y^2+xy=x^3+1+\epsilon$ and keep the ordinary translation factor fixed. Its nonidentity two-torsion component is $\operatorname{Spec}R[a]/(a^2-\epsilon)$, obtained by putting $x=0,y=1+a$. It is finite flat despite lacking an $R$-point. Its absolute tangent dimension is one smaller than that of the neutral $\mu_2$ component. Thus the nonsmooth-torsor behavior used to diagnose the published inference occurs in an Igusa family itself.
